\chapter{New numbers for old recurrences}
\label{app:finite-fields}
\chaptertag{An optional later route: extensions, Frobenius, and Fibonacci periods}

\begin{bigidea}
Chapter~\ref{chap:modular} proves Fibonacci periodicity using only
reversible pairs of residues. This appendix answers the later question:
how can the roots of $X^2-X-1$ help compute a period when those roots
are not in $\F_p$? First we build a field containing them and learn its
arithmetic. Only then do we use Frobenius and multiplicative orders.
\end{bigidea}

Read this after modular inverses and interpolation
(Chapter~\ref{chap:modular}), Fermat's theorem
(Chapter~\ref{chap:necklaces}), orders in finite groups
(Chapter~\ref{chap:groups}), and the prime binomial identity
(Chapter~\ref{chap:binomial}). The current modular-arithmetic course can
stop at the reversible-pair proof and return here later.

\section{Build an extension before taking powers in it}
\label{sec:quadratic-field-arithmetic}

The squares modulo $7$ are $0,1,2,4$. There is no element of $\F_7$
whose square is $5$, so writing $\sqrt5$ does not yet give us a number
we can calculate with in that field.

Introduce a symbol $s$ with the rule $s^2=5$, and use pairs
\[
 a+bs,\qquad a,b\in\F_7.
\]
Addition is coefficientwise. Multiply by distributing, then replace
$s^2$ by $5$:
\[
 (a+bs)(c+es)=(ac+5be)+(ae+bc)s.
\]
All coefficients are reduced modulo $7$. For example,
\[
 (2+s)(3+2s)=16+7s=2.
\]
Formally, this is polynomial arithmetic modulo $X^2-5$:
\[
 K=\F_7[X]/(X^2-5).
\]
The notation means that two polynomials represent the same element when
their difference is a multiple of $X^2-5$. Polynomial division gives a
unique representative $a+bX$, so there are exactly $49$ elements.
The usual arithmetic laws descend from polynomial addition and
multiplication, just as they did for integer residues.

\begin{proposition}[Division in the new arithmetic]
Every nonzero $a+bs$ has an inverse, given by
\[
 (a+bs)^{-1}=\frac{a-bs}{a^2-5b^2}.
\]
Consequently $K$ is a field of $49$ elements, denoted $\F_{49}$.
\end{proposition}
\begin{proof}
The product of $a+bs$ and $a-bs$ is $a^2-5b^2\in\F_7$.
This denominator is nonzero for a nonzero pair: if $b=0$, it is $a^2$;
if $b\ne0$ and it vanished, then $(a/b)^2=5$, contradicting the list
of squares. Its inverse exists in $\F_7$ and can be computed by
extended Euclid. Multiplying gives the claimed inverse.
\end{proof}

For instance, $4-5=-1=6$ has inverse $6$ modulo $7$, so
\[
 (2+s)^{-1}=6(2-s)=5+s.
\]
Indeed, $(2+s)(5+s)=15+7s=1$.
Extended gcd has not disappeared: it now performs the division of
coefficients inside the extension.

\begin{warningbox}
$\F_{49}$ is not $\Z/49\Z$. In $\F_{49}$, seven copies of $1$ sum to
$0$, and every nonzero element is invertible. In $\Z/49\Z$, the nonzero
residue $7$ satisfies $7\cdot7=0$ and has no inverse.
\end{warningbox}

The construction works for any odd prime $p$ and any nonsquare
$d\in\F_p$: the pairs $a+bs$ with $s^2=d$ form a field of $p^2$
elements. The same inverse formula uses $a^2-db^2$ in the denominator.
Equivalently, $X^2-d$ has no root in $\F_p$, so, being quadratic,
it is irreducible. If a chosen polynomial factors, its quotient need
not be a field; irreducibility is doing the work.

\section{Frobenius moves the new elements}
\label{sec:frobenius}

A field automorphism is a bijection from a field to itself preserving
addition, multiplication, and $1$. On a finite field $K$ containing
$\F_p$, define
\[
 \operatorname{Fr}(u)=u^p.
\]

\begin{theorem}[Frobenius on a finite extension]
The map $\operatorname{Fr}$ is an automorphism of $K$. Its fixed
elements are exactly the embedded prime field $\F_p$.
\end{theorem}
\begin{proof}
The coefficient argument in Chapter~\ref{chap:binomial} gives
$(u+v)^p=u^p+v^p$ in $K$, and $(uv)^p=u^pv^p$ follows from
commutativity. Also $1^p=1$. If $u^p=v^p$, then $(u-v)^p=0$;
a nonzero field element cannot have a zero power, so $u=v$.
Thus the map is injective, and finiteness makes it surjective.

Fermat's theorem fixes every element of $\F_p$. There can be no
additional fixed elements: the nonzero polynomial $X^p-X$ has degree
$p$, so the root bound from interpolation allows at most $p$ roots in
$K$. The prime-field elements already supply all $p$ of them.
\end{proof}

\begin{example}[Frobenius is conjugation in $\F_{49}$]
Since $s^2=5$,
\[
 s^7=s(s^2)^3=5^3s=125s=-s.
\]
For $a,b\in\F_7$, the sum identity and Fermat now give
\[
 (a+bs)^7=a^7+b^7s^7=a-bs.
\]
Thus Frobenius exchanges $s$ and $-s$, fixes the $7$ constant elements,
and applying it twice returns every element to itself. This is a
nontrivial symmetry of the field, although it acts as the identity on
the coefficients.
\end{example}

More generally, if $s^2=d$ is nonsquare over $\F_p$ with $p$ odd,
then $(s^p)^2=d^p=d$, so $s^p$ is either $s$ or $-s$.
It cannot be $s$, since $s$ is not in the fixed prime field.
Therefore $(a+bs)^p=a-bs$ in this quadratic extension as well.
Finiteness was essential to the surjectivity argument; no assertion
about arbitrary infinite fields is needed here.

\section{Return to the Fibonacci roots}
\label{sec:fibonacci-field-periods}

Let $p\ne2,5$ be prime. If $5$ is a square in $\F_p$, work in $K=\F_p$;
otherwise construct $K=\F_p(s)$ with $s^2=5$ as above.
In either case, $K$ contains the two distinct nonzero roots
\[
 \alpha=\frac{1+s}{2},\qquad\beta=\frac{1-s}{2}
\]
of $X^2-X-1$. In particular,
$\alpha+\beta=1$, $\alpha\beta=-1$, and $\alpha-\beta=s\ne0$.

The formula
\[
 F_n=\frac{\alpha^n-\beta^n}{\alpha-\beta}\qquad\text{in }K
\]
now makes sense: all its divisions are by invertible field elements.
It is proved by checking $n=0,1$ and observing that each of
$\alpha^n,\beta^n$ satisfies the Fibonacci recurrence.
The resulting values are the images of the ordinary Fibonacci numbers
in the subfield $\F_p$.

\begin{proposition}[A period is a simultaneous multiplicative order]
For $p\ne2,5$,
\[
 \pi(p)=\operatorname{lcm}\bigl(\operatorname{ord}(\alpha),
                                  \operatorname{ord}(\beta)\bigr).
\]
\end{proposition}
\begin{proof}
A return $F_t=0$ first says $\alpha^t=\beta^t=c$.
Then $F_{t+1}=c(\alpha-\beta)/(\alpha-\beta)=c$, so the full
return condition $(F_t,F_{t+1})=(0,1)$ is equivalent to
$\alpha^t=\beta^t=1$. The least such positive $t$ is the displayed
least common multiple.
\end{proof}

\subsection*{When the roots are already in the prime field}

If $5$ is a square in $\F_p$, Fermat gives
$\alpha^{p-1}=\beta^{p-1}=1$, and hence
\[
 \pi(p)\divides p-1.
\]
For $p=11$, choose $s=4$. The roots are $\alpha=8$, $\beta=4$.
Their orders are $10$ and $5$: $8^5=-1$, $8^2=9\ne1$, while
$4^5=1$ and $4\ne1$. Thus $\pi(11)=10$.
The ambient orders divide $10$ by Fermat, which certifies the minimality
claims in these power checks.

\subsection*{When Frobenius exchanges the roots}

If $5$ is nonsquare, Frobenius takes $s$ to $-s$ and therefore
\[
 \alpha^p=\beta,\qquad\beta^p=\alpha.
\]
Multiplying by the original root gives the sharper information
\[
 \alpha^{p+1}=\beta^{p+1}=\alpha\beta=-1.
\]
Consequently
\[
 \pi(p)\divides2(p+1),\qquad \pi(p)\ndivides p+1.
\]
The minus sign cannot be dropped: $p$ is odd.
Indeed, $F_{p+1}=0$ but $F_{p+2}=-1$, so the state at $p+1$ is
$(0,-1)$, not a restart at $(0,1)$.

\begin{example}[The period modulo $7$]
In the field constructed earlier, $\alpha=4+4s$ and $\beta=4-4s$.
Repeated squaring gives
\[
 \alpha^2=5+4s,\qquad \alpha^4=5s,\qquad
 \alpha^8=-1,\qquad\alpha^{16}=1.
\]
Thus $\alpha$ has order $16$: every proper divisor of $16$ divides $8$,
and $\alpha^8\ne1$. Both roots have $16$th power $1$, so
$\pi(7)=16$. This computation uses the arithmetic of the extension
explicitly; it does not treat the missing square root as an ordinary
residue modulo $7$.
\end{example}

\section{From a bound to the exact period}
\label{sec:exact-fibonacci-period}

The divisibilities above supply a candidate multiple, not always the
least period. For instance, $\pi(29)=14$, a proper divisor of $28$.

We can decide which case applies without listing all squares.
There are $(p-1)/2$ nonzero squares, since $x^2=y^2$ exactly when
$x=\pm y$. Fermat makes each such square a root of
$X^{(p-1)/2}-1$, and the root bound allows no others.
Thus $5$ is a square exactly when $5^{(p-1)/2}=1\pmod p$.
Repeated squaring computes this test efficiently.

A practical procedure after finding a valid multiple $B$ is:
\begin{enumerate}
  \item Factor $B$ and set $t=B$.
  \item For each prime divisor $q$ of $t$, test whether
  $(F_{t/q},F_{t/q+1})=(0,1)\pmod p$.
  If so, replace $t$ by $t/q$ and continue removing factors.
  \item Stop when no prime factor can be removed. The remaining $t$
  is $\pi(p)$.
\end{enumerate}
To justify the stopping rule, $\pi(p)$ divides the current $t$ at every
stage. If $t>\pi(p)$, some prime divides $t/\pi(p)$, and removing that
prime still leaves a multiple of the true period. Thus a larger value
cannot survive all the tests.

These tests can be made without computing all preceding Fibonacci
numbers. The doubling identities are
\[
 F_{2n}=F_n(2F_{n+1}-F_n),\qquad
 F_{2n+1}=F_n^2+F_{n+1}^2.
\]
For a proof, the recurrence gives
\[
 \begin{pmatrix}0&1\\1&1\end{pmatrix}^{n}
 =\begin{pmatrix}F_{n+1}-F_n&F_n\\F_n&F_{n+1}\end{pmatrix}.
\]
Square this matrix and compare its upper-right and lower-right entries.
Thus a pair at index $n$ gives the pair at $2n$; one additional recurrence
step gives the pair at $2n+1$. Reading the binary digits of an index
computes its pair by repeated doubling and reduction.

For $p=29$, $11^2=5$ in $\F_{29}$, so start with $B=28$.
The pair at $14$ is $(0,1)$, reducing the candidate to $14$.
The pairs at $7$ and $2$ are $(13,21)$ and $(1,2)$, respectively,
so neither prime factor of $14$ can be removed. Hence $\pi(29)=14$.

\subsection*{The exceptional primes}

For $p=2$, direct pair iteration gives $\pi(2)=3$.
For $p=5$, the characteristic polynomial has a repeated root, so the
distinct-root argument does not apply. Put
$A=\begin{psmallmatrix}0&1\\1&1\end{psmallmatrix}$. Direct computation
gives $A^5=3I\pmod5$, hence $A^{20}=I$.
But $A^{10}=4I\ne I$, and $A^4$ has off-diagonal entry $F_4=3\ne0$.
Every proper divisor of $20$ divides $10$ or $4$, so the order of $A$
is $20$. The displayed matrix formula shows that $A^t=I$ exactly when
$(F_t,F_{t+1})=(0,1)$; therefore $\pi(5)=20$.

\begin{checkpoint}
Keep three questions separate. Reversible finite states prove that a
period exists. Extension-field arithmetic and Frobenius give useful
multiples of that period. Order tests identify the least period.
For small moduli, simply following the pair orbit remains a complete
method and requires none of the later field theory.
\end{checkpoint}

\begin{exercisebox}
\begin{enumerate}
  \item In $\F_7(s)$ with $s^2=5$, compute $(1+2s)(3+s)$ and
  $(1+2s)^{-1}$. Check the inverse by multiplication.
  \exlevel{Fluency}
  \item Construct $\F_4$ using $t^2=t+1$ over $\F_2$.
  List its four elements and invert each nonzero element.
  Compute the square map and identify its fixed elements.
  \exlevel{Reasoning}
  \item In $\F_7(s)$, show that $u^7=u$ exactly when $u\in\F_7$,
  and that $u^{49}=u$ for every element. Explain which calculation
  exhibits Frobenius as a nonidentity automorphism.
  \exlevel{Reasoning}
  \item In $\F_{29}$, verify that the roots of $X^2-X-1$ are $6$ and
  $24$, of orders $14$ and $7$. Deduce the period by their orders
  and compare with the factor-removal calculation.
  \exlevel{Reasoning}
  \item When $5$ is nonsquare in $\F_p$ for odd $p$, prove directly
  that the Fibonacci pair at $p+1$ is $(0,-1)$. Use $p=3$ to connect
  the result with the pair table in Chapter~\ref{chap:modular}.
  \exlevel{Synthesis}
  \item Implement the doubling identities and compare the exact
  period procedure with direct pair iteration for several small
  primes. Treat $2$ and $5$ separately, and report candidate multiples
  as well as the least periods. \exlevel{Exploration}
\end{enumerate}
\end{exercisebox}
